\begin{table}[htbp]
\centering
\caption{How far simulated $p$-hacking moves the $t$-curve from the honest set}
\label{tab:brodeur_all_power_dw}
\footnotesize
\begin{tabular}{lcc}
\toprule
$\tau$ & Maximum of 8 & Threshold \\
\midrule
0.0 & 0.000 & 0.000 \\
0.1 & 0.005 & 0.168 \\
0.2 & 0.051 & 0.455 \\
0.3 & 0.150 & 0.784 \\
0.4 & 0.293 & 1.138 \\
0.5 & 0.484 & 1.509 \\
0.6 & 0.747 & 1.894 \\
0.7 & 1.125 & 2.289 \\
0.8 & 1.788 & 2.694 \\
0.9 & 3.496 & 3.107 \\
1.0 & 7.289 & 3.528 \\
\midrule
\multicolumn{3}{l}{\textit{Share $\tau$ at which $d_w$ equals the application's $\hat B$}} \\
RCT ($\hat B = 2.8$) & 0.86 & 0.83 \\
IV ($\hat B = 3.1$) & 0.88 & 0.90 \\
DID ($\hat B = 0.9$) & 0.64 & 0.33 \\
RDD ($\hat B = 1.3$) & 0.73 & 0.44 \\
\bottomrule
\end{tabular}
\par\smallskip
\begin{minipage}{\linewidth}\footnotesize\textit{Notes:} Entries: $d_w \times 10^3$, the weighted distance between the $t$-curve produced when a share $\tau$ of studies $p$-hack and the set of $t$-curves possible without selective reporting, computed exactly rather than simulated, with the projection test's coefficients ($J = 40$, $J_0 = 20$, symmetrized and shifted by 1.96) and its grid of 3{,}000 points on $[-6.5, 6.5]$ plus 9999. Effects are drawn from a Gamma with shape 1.413 and rate 0.770, as in the simulations, with no rounding. Maximum of 8: the largest of 8 $t$-statistics; Threshold: a second draw if $t \le 1.96$. The breakdown statistic $\hat B$ is a lower 95\% confidence bound on the same distance for the reported $t$-curve, so a design whose $\hat B$ exceeds $d_w$ at $\tau$ is at least as far from honest reporting as that form of $p$-hacking by a share $\tau$ of its studies would make it. The shares depend on the distribution of effects and the form of $p$-hacking; they convey magnitudes, not estimates of the extent of $p$-hacking.\end{minipage}
\end{table}
